\documentclass[10pt,a4paper]{amsart}
\usepackage[margin=19mm]{geometry}
\usepackage{amsmath,amssymb,mathtools}
\usepackage[scr=rsfs,cal=euler]{mathalfa}
\usepackage{xfrac}
\usepackage[unicode,hidelinks,bookmarks=false]{hyperref}
\newcommand{\ie}{i.e.\ }
\newcommand{\cf}{cf.\ }
\newcommand{\resp}{respectively\ }
\newcommand{\Subsection}{\subsection}
\providecommand{\nobreakdash}{\nobreak\hbox{-}}
\setlength{\emergencystretch}{2em}
\multlinegap0pt
\usepackage[shortlabels]{enumitem}
\usepackage{amssymb}
\usepackage{mathtools}
\usepackage{tikz}
\newtheorem{theorem}{Theorem}
\title{Time Structure in Infinite Extensive Games}
\author{Shravan Luckraz}
\date{Source: arXiv:2606.26582v3 (CC BY 4.0)}
\begin{document}
\maketitle

\noindent\emph{Source and licence.} Time Structure in Infinite Extensive Games. Version arXiv:2606.26582v3 (CC BY 4.0), distributed by arXiv under the Creative Commons Attribution 4.0 International licence, http://creativecommons.org/licenses/by/4.0/, as recorded in the arXiv licence metadata for this exact version and shown on the abstract page https://arxiv.org/abs/2606.26582v3. Full licence notice: \texttt{licenses/arxiv-2606.26582-CC-BY-4.0.txt}.\par\smallskip

\noindent\textit{Editorial scope.} Counted unit: the article's theorem theorem (source0.tex lines 1348--1358). The article's complete original proof of it is delivered with it (source0.tex lines 1360--1401, one \texttt{proof} environment of 42 lines). No mathematics is written, repaired or re-derived: the statement and the proof are the article's own lines, in the article's own notation, and no retained line is substituted or re-flowed.  No further source-local context is needed: every cross-reference of every retained line resolves inside the two delivered spans.  Nothing was omitted from any retained span, so the package carries no omission record.  No retained line contains a citation, so the wrapper carries no bibliography and no bibliography entry.  No retained line contains a figure environment, an image-inclusion command or drawing code, so no image is delivered and the package declares no asset dependency.  Editorial formatting only, in addition: an AMS \texttt{amsart} preamble that restates the article's own declarations and macros used by the retained lines, namely the environments \texttt{theorem} on separate counters, together with the standard packages those lines need.  The retained environments print in this document as Theorem 1 in order of appearance, whereas the article prints the same items with the numbering of its own full document.  The complete native source of the article as distributed in the arXiv e-print for this exact version is bundled unchanged as \texttt{sources/203699/source0.tex}.
\begin{theorem}
Let $V$ be a second‑countable metric space and $\prec\subseteq V\times V$
a strict partial order satisfying:
\begin{enumerate}
  \item $\prec$ is acyclic,
  \item $\prec$ is topologically closed in $V\times V$,
  \item there exists a countable order‑dense subset $D\subset V$.
\end{enumerate}
Then there exists a continuous $f:V\to\mathbb R$ such that
$x\prec y\Rightarrow f(x)<f(y)$.
\end{theorem}

\begin{proof}
The proof adapts Debreu's separability argument for preorders to strict
partial orders.

Let $D=\{d_1,d_2,\dots\}$ be order‑dense. For each $d_i$ define


\[
L_i:=\{x\in V: x\prec d_i\},\qquad R_i:=\{x\in V: d_i\prec x\}.
\]


Because the graph of $\prec$ is closed, the sets $L_i$ and $R_i$ are
open: if $x\prec d_i$ then by closedness there exists a neighborhood $U$
of $x$ such that for all $u\in U$ we still have $u\prec d_i$. Thus $L_i$
is open; similarly for $R_i$.

Enumerate rationals $\mathbb Q=\{q_1,q_2,\dots\}$. We will construct open
sets $U_{q_k}$ recursively so that:
\begin{enumerate}[label=(\roman*)]
  \item $U_{q_k}\subseteq U_{q_{k'}}$ whenever $q_k<q_{k'}$,
  \item for every $x\prec y$ there exist rationals $q<q'$ with $x\in U_q$
    and $y\notin\overline{U_q}$.
\end{enumerate}

Initialize $U_{q_1}=\varnothing$. At stage $k$, suppose $U_{q_j}$ has been
defined for $j<k$ satisfying monotonicity. Consider the set of ordered
pairs $(x,y)$ not yet separated by earlier $U_{q_j}$. For each such pair
pick $d_i\in D$ with $x\prec d_i\prec y$ (order‑density). Using second
countability, choose for each such $d_i$ a small open neighborhood $W_i$
contained in $R_i$ and disjoint from the closure of the current
$U_{q_{k-1}}$. Let $U_{q_k}$ be the union of $U_{q_{k-1}}$ with finitely
many such neighborhoods $W_i$ chosen so that the new $U_{q_k}$ separates
a finite but increasing set of previously unseparated pairs. Because the
rationals are countable and second countability allows us to refine
neighborhoods, this recursive construction can be carried out so that
every ordered pair is eventually separated at some stage.

Define $f(x)=\inf\{q\in\mathbb Q: x\in U_q\}$. By construction $f(x)<f(y)$
whenever $x\prec y$. Standard refinements of the construction (shrinking
neighborhoods and using second countability) ensure $f$ is continuous.
\end{proof}

\end{document}
